\begin{abstract}
We present \repo{} as a software-and-methods workflow for nonlinear finite-element
solves and optimization that spans \purejax{} references, \fenics{} references
where they exist, and a maintained distributed \jaxpetsc{} implementation.
Across the maintained benchmark suite we cover nonlinear scalar problems
($p$-Laplace and Ginzburg--Landau), large-deformation hyperelasticity,
two-dimensional and three-dimensional Mohr--Coulomb plasticity, and topology
optimization. The methodological focus is a solver-centric combination of JAX
for local differentiation with PETSc for sparse assembly, globalization,
Krylov solves, and preconditioning, together with three second-order routes:
element automatic differentiation, constitutive automatic differentiation, and
colored sparse finite-difference recovery.

The evidence is separated into validation and performance layers. For the
three-dimensional heterogeneous Plasticity3D benchmark, a same-case
source-faithfulness comparison against the source-style endpoint surrogate is
tight, with relative differences
\num{3.877e-5} in work, \num{3.517e-3} in displacement, and
\num{8.720e-3} in deviatoric-strain norm. A separate fixed-lambda
source-operator diagnostic uses the matched glued-bottom boundary contract and
passes the configured field-agreement thresholds, with
$u_{\max}(\lambda)$ relative $L^2$ difference \num{6.396e-6} and endpoint
displacement relative $L^2$ difference \num{8.299e-4}. This validates the
implemented endpoint surrogate within the tested fixed-load contract, but not
as a path-consistent replacement for an incremental elastoplastic history
solve. On the locked $P_4(L_1)$,
$\lambda_{\mathrm{sr}}=1.5$, 8-rank derivative ablation, constitutive AD has
the lowest median wall time among the three tested routes while reaching the
same terminal state as element AD and colored SFD. A serial hyperelastic
companion comparison against JAX-FEM on the same mesh and displacement schedule
passes a 5\% fairness gate with repository-energy and field differences below
$10^{-3}$; this evidence is hyperelastic-only.

The largest maintained distributed timing study remains the promoted
\plasticitythree{} study on $P_4(L_2)$ with \localpmg{}, which attains
$\num{18.05}\times$ end-to-end speedup on 32 MPI
ranks relative to one rank at the same promoted nonlinear target. The topology
optimization study complements this with a distributed design-and-mechanics
workflow that remains operational at $768\times 384$ resolution and 32 ranks.
Within that scoped evidence, the paper's contribution is a maintained
implementation pattern: derivative-rich local models, sparse nonlinear solver
policy, and distributed execution can be kept together in one reproducible
workflow across a heterogeneous family of FEM problems.
\end{abstract}
